% ============================================================
%  Suwoong Yeom — Curriculum Vitae
%  pdfLaTeX / XeLaTeX / LuaLaTeX 어느 엔진으로도 컴파일됩니다 (Overleaf 기본값 OK).
%
%  자주 쓰는 명령
%    \me          본인 이름 (굵게)
%    \eq          equal contribution 표시 (*)
%    \corr        교신저자 표시 (†)
%    \verify{…}   확인이 필요한 문장 (빨간 글씨) — 내용 확인 후 \verify{ } 를 벗겨내세요
%    \entry{제목}{장소}{소속}{기간}
%    \pub{번호}{저자}{제목}{학회/저널}{비고 — 없으면 {} 로 비움}
% ============================================================
\documentclass[10pt]{article}

\usepackage{iftex}
\ifPDFTeX
  \usepackage[T1]{fontenc}
  \usepackage[utf8]{inputenc}
  \usepackage{lmodern}
\else
  \usepackage{fontspec}
\fi

\usepackage[a4paper, top=1.8cm, bottom=2.2cm, left=1.9cm, right=1.9cm]{geometry}
\usepackage{titlesec}
\usepackage{enumitem}
\usepackage{xcolor}
\usepackage{fancyhdr}
\usepackage{lastpage}
\usepackage[hidelinks]{hyperref}
\hypersetup{pdftitle={Suwoong Yeom - Curriculum Vitae}, pdfauthor={Suwoong Yeom}}

% ---------- 페이지 ----------
\pagestyle{fancy}
\fancyhf{}
\renewcommand{\headrulewidth}{0pt}
\fancyfoot[L]{\small Last modified: \today}
\fancyfoot[C]{\small \thepage\ of \pageref{LastPage}}

\setlength{\parindent}{0pt}
\raggedbottom

% ---------- 섹션 제목: 작은 대문자 + 밑줄 ----------
\titleformat{\section}{\large\scshape\raggedright}{}{0em}{}[\vspace{-0.3ex}\titlerule]
\titlespacing*{\section}{0pt}{1.6ex}{1.0ex}
\titleformat{\subsection}{\normalsize\raggedright}{}{0em}{}
\titlespacing*{\subsection}{0pt}{0.8ex}{0.2ex}

% ---------- 목록 ----------
\newlist{entries}{itemize}{1}
\setlist[entries]{label={}, leftmargin=1em, itemsep=1.0ex, topsep=0.2ex, parsep=0pt}
\newlist{details}{itemize}{1}
\setlist[details]{label=\textbullet, leftmargin=1.6em, itemsep=0.1ex, topsep=0.2ex, parsep=0pt}
\newlist{pubs}{itemize}{1}
\setlist[pubs]{leftmargin=3.4em, labelwidth=2.8em, labelsep=0.6em, align=right,
               itemsep=1.0ex, topsep=0.4ex, parsep=0pt,
               before={\interlinepenalty=10000 \clubpenalty=10000 \widowpenalty=10000}}  % 논문 한 편이 두 쪽에 걸쳐 잘리지 않게

% ---------- 명령 ----------
\newcommand{\me}{\textbf{Suwoong Yeom}}
\newcommand{\eq}{\textsuperscript{*}}
\newcommand{\corr}{\textsuperscript{\dag}}
\newcommand{\verify}[1]{\textcolor{red}{#1}}
% 전부 확인했으면 위 줄 대신 아래 줄을 쓰면 빨간 글씨가 한 번에 검정으로 바뀝니다.
% \newcommand{\verify}[1]{#1}

% \entry{제목}{장소}{소속}{기간}
\newcommand{\entry}[4]{%
  \item \textbf{#1}\hfill #2\\*
  \textit{#3}\hfill\textit{#4}}

% \project{과제명}{기관 · 역할}{기간}
\newcommand{\project}[3]{%
  \item \textit{#2}\hfill\textit{#3}\\*
  {\hyphenpenalty=10000 \exhyphenpenalty=10000 \raggedright\textbf{#1}\par}}  % 과제명은 단어 중간에서 끊지 않음

% \pub{번호}{저자}{제목}{학회/저널}{비고}
\newcommand{\pub}[5]{%
  \item[{[#1]}] #2\\*
  \textit{#3}\\*
  #4%
  \if\relax\detokenize{#5}\relax\else\\* #5\fi}

\begin{document}

% ============================================================
%  이름 / 연락처
% ============================================================
\begin{center}
  {\fontsize{26}{30}\selectfont\bfseries Suwoong Yeom}\\[5pt]
  {\small
    \href{tel:+821088599561}{+82-10-8859-9561} \textbar{}
    \href{mailto:suwoong96@sogang.ac.kr}{suwoong96@sogang.ac.kr} \textbar{}
    \href{https://yeomsuwoong.github.io}{yeomsuwoong.github.io} \textbar{}
    \href{https://scholar.google.com/citations?hl=en&user=7wG0GVcAAAAJ}{Google Scholar}%
    % 필요하면 아래 줄 주석을 풀어 추가하세요.
    % \textbar{} \href{https://github.com/YeomSuWoong}{GitHub}%
    % \textbar{} \href{https://www.linkedin.com/in/suwoong-yeom-825599251/}{LinkedIn}%
    \\[1pt]
    VDSLab, Department of Electronic Engineering, Sogang University, Seoul, Republic of Korea}
\end{center}

% ============================================================
\section{Research Interests}
% ============================================================
\begin{entries}
  \item 3D/4D reconstruction and neural rendering with NeRF and Gaussian Splatting;
  spatial understanding in feed-forward models; world models for digital twins and physical AI.
\end{entries}

% ============================================================
\section{Education}
% ============================================================
\begin{entries}
  \entry{Integrated M.S./Ph.D. in Electronic Engineering}{Seoul, Republic of Korea}
        {Sogang University (VDSLab)}{2023.03 -- Present}
  \begin{details}
    \item Overall GPA: 4.15/4.5
    \item Advisor: Prof. Suk-Ju Kang
  \end{details}

  \entry{B.S. in Mechanical Engineering}{Suwon, Republic of Korea}
        {Ajou University}{2016.03 -- 2022.08}
  \begin{details}
    \item Overall GPA: 3.42/4.5
  \end{details}
\end{entries}

% ============================================================
\section{Research Experience}
% ============================================================
\begin{entries}
  \entry{Visiting Researcher}{West Lafayette, IN, USA}
        {NEISLab, Purdue University}{2026.06 -- 2026.07}
  \begin{details}
    \item Advisor: Prof. Younghyun Kim
  \end{details}
\end{entries}

% ============================================================
\section{Publications}
% ============================================================
{\small \eq Equal contribution (co-first authors) \quad \corr Corresponding author}

\subsection{International Conference}
\begin{pubs}
  \pub{C2}
      {\me\eq, Joonsik Nam\eq, Seunggyu Choi\eq, Lucas Yunkyu Lee, Sangmin Kim, Jaesik Park,
       Joonsoo Kim, Kug-jin Yun, Kyeongbo Kong\corr, and Suk-Ju Kang\corr}
      {TRiGS: Temporal Rigid-Body Motion for Scalable 4D Gaussian Splatting}
      {European Conf. on Computer Vision (\textbf{ECCV}), 2026}
      {}

  \pub{C1}
      {\me, Hosung Son, Chanhee Kang, Joonsoo Kim, Kug-jin Yun, Won-Sik Cheong, and Suk-Ju Kang}
      {Motion Mask-driven Improvements in Monocular Dynamic Novel View Synthesis}
      {Int. Conf. on Electronics, Information, and Communications (\textbf{ICEIC}), 2024}
      {(\textbf{Oral} presentation)}
\end{pubs}

\subsection{International Journal}
\begin{pubs}
  \pub{J1}
      {\me, Hosung Son, Chanhee Kang, Eunho Shin, Joonsoo Kim, Kug-jin Yun, and Suk-Ju Kang}
      {ICSD-NeRF: Independent Canonical Spaces for Enhanced Dynamic Scene Modeling in Neural Radiance Fields}
      {IEEE Transactions on Computational Imaging (\textbf{TCI}), vol.~12, pp.~242--255, 2026}
      {}
\end{pubs}

\subsection{Workshop}
\begin{pubs}
  \pub{W3}
      {Seunggyu Choi\eq, \me\eq, and Suk-Ju Kang\corr}
      {The Decoder Already Knows Where to Densify: Attention-Gated Capacity Allocation for Feed-Forward 3D Gaussian Splatting}
      {3DWM Workshop, European Conf. on Computer Vision (\textbf{ECCV}), 2026}
      {}

  \pub{W2}
      {Chaewon Lee\eq, Seunggyu Choi\eq, \me\eq, Joonsoo Kim, Kugjin Yun, and Suk-Ju Kang\corr}
      {Seamless Spatial Extension in 3D Gaussian Splatting via Auxiliary Wild Views}
      {TwinWorld Workshop, European Conf. on Computer Vision (\textbf{ECCV}), 2026}
      {}

  \pub{W1}
      {Jeongseon Oh\eq, \me\eq, ChangHee Yang\eq, and Suk-Ju Kang\corr}
      {Dual-PAG: Bidirectional Perturbed-Attention Guidance for Residue-Free Video Editing}
      {F2S Workshop, Int. Conf. on Machine Learning (\textbf{ICML}), 2026}
      {}
\end{pubs}

\subsection{Under Review}
\begin{pubs}
  \pub{U4}
      {\me, Geonho Cha, Dongyoon Wee,
       Kyeongbo Kong\corr, and Suk-Ju Kang\corr}
      {Learning to Remember and Evolve: Object State Memory for Video World Models}
      {International Conference on Learning Representation (\textbf{ICLR})}
      {}
  \pub{U3}
      {Seunggyu Choi\eq, \me\eq, Joonsik Nam, Geonho Cha, Dongyoon Wee,
       and Suk-Ju Kang\corr}
      {Learning to Optimize: Prior-Preserving Point-to-Gaussian Lifting for Feed-Forward Novel View Synthesis}
      {The Association for the Advancement of Artificial Intelligence (\textbf{AAAI})}
      {}
  \pub{U2}
      {\me\eq, Jimin Roh\eq, Eunho Shin\eq, Songju Na, Joonsoo Kim, Kug-jin Yun,
       Kyeongbo Kong\corr, and Suk-Ju Kang\corr}
      {EMDGS: Enhancing Monocular Dynamic Gaussian Splatting via Adaptive Pipeline Integration}
      {IEEE Transactions on Circuits and Systems for Video Technology (\textbf{TCSVT})}
      {}

  \pub{U1}
      {Taewoo Kim\eq, \me\eq, Jaehyun Pyun\eq, Geonho Cha, Dongyoon Wee, Joonsik Nam,
       Yun-Seong Jeong, Kyeongbo Kong\corr, and Suk-Ju Kang\corr}
      {HOIGS: Human-Object Interaction Gaussian Splatting}
      {IEEE Transactions on Multimedia (\textbf{TMM})}
      {}
\end{pubs}

% ============================================================
\section{Research Projects}
% ============================================================
%  \project{과제명}{기관}{기간} — 각 항목 위 주석이 한글 원제입니다.
%  수행 내용을 적고 싶으면 과제 아래에 \begin{details} \item … \end{details} 를 추가하세요.
%  빨간 글씨(관련 논문)는 초안이니 확인 후 고쳐 주세요.
\begin{entries}
  % 보조 영상 기반 동적 장면 공간 확장 방안 연구
  \project{Spatial Extension of Dynamic Scenes Using Auxiliary Views}
          {Electronics and Telecommunications Research Institute (ETRI)}{2026.01 -- Present}

  % 시공간 객체 상태 일관성 유지 가능 비디오 월드 모델 개발
  \project{Video World Models with Spatiotemporally Consistent Object States}
          {NAVER Cloud}{2026.01 -- Present}
  \begin{details}
    \item Related publications: [U4], [U3]
  \end{details}

  % 인간 객체 상호작용 가우시안 스플래팅 기술 개발
  \project{Human-Object Interaction Gaussian Splatting}
          {NAVER Cloud}{2025.06 -- 2025.12}
  \begin{details}
    \item Related publication: [U1]
  \end{details}

  % 보조 영상 기반 정적 장면 공간 확장 및 동적 장면 표현 고도화 연구
  \project{Spatial Extension of Static Scenes Using Auxiliary Views and Enhanced Dynamic Scene Representation}
          {Electronics and Telecommunications Research Institute (ETRI)}{2025.01 -- 2025.12}
  \begin{details}
    \item Related publications: [C2], [W2]
  \end{details}

  % 생성모델 기반 3차원 공간 표현 인페인팅 알고리즘 연구
  \project{Generative Model-Based Inpainting for 3D Scene Representations}
          {Electronics and Telecommunications Research Institute (ETRI)}{2024.01 -- 2024.12}

  % 실사 영상 3차원 공간정보 획득 방안 연구
  \project{Acquiring 3D Spatial Information from Real-World Videos}
          {Electronics and Telecommunications Research Institute (ETRI)}{2023.01 -- 2023.12}
  \begin{details}
    \item Related publications: [C1], [J1]
  \end{details}
\end{entries}

% ============================================================
\section{Mentoring}
% ============================================================
\begin{details}[leftmargin=2.6em]
  \item \textbf{Seunggyu Choi}, M.S. Student --- co-first-authored [W3], ECCV 2026 3DWM Workshop
  \item \textbf{Chaewon Lee}, Undergraduate Researcher --- co-first-authored [W2], ECCV 2026 TwinWorld Workshop
  \item \textbf{Jeongseon Oh}, Integrated M.S./Ph.D. Student --- co-first-authored [W1], ICML 2026 F2S Workshop
\end{details}

% ============================================================
\section{Languages}
% ============================================================
\begin{details}[leftmargin=2.6em]
  \item \textbf{Korean}: Native
  \item \textbf{English}: OPIc IH (Intermediate High)
\end{details}

% ============================================================
%  아래 섹션은 실제로 해당하는 것만 주석(%)을 풀고 채우세요.
%  해당 사항이 없는 섹션은 그대로 두면 PDF에 나오지 않습니다.
% ============================================================

% \section{Technical Skills}
% \begin{details}
%   \item \textbf{Programming}: Python, PyTorch, CUDA, C++
%   \item \textbf{3D/Graphics}: COLMAP, Nerfstudio, gsplat, Blender
% \end{details}

% \section{Honors and Awards}
% \begin{details}
%   \item \textbf{Award Name}, Organization, Mon.\ YYYY
% \end{details}

% \section{Academic Service}
% \begin{details}
%   \item Reviewer: CVPR, ICCV, ECCV, NeurIPS, ...
% \end{details}

\end{document}
